\section{Comparative Intervention Map}

\begin{itemize}[nosep]
  \item \textbf{Window capture.} Target: timing and entry leverage.
        Evidence lane: strong + design inference.
        Use: move the attempt to a boundary such as wake-up, post-break
        re-entry, or commute transition.
  \item \textbf{State preloading.} Target: switching cost and path clarity.
        Evidence lane: medium + design inference.
        Use: prepare the exact next page, sentence, route, or tool before the
        window opens.
  \item \textbf{Environmental re-weighting.} Target: cue landscape and local
        reinforcement.
        Evidence lane: strong/medium + design inference.
        Use: remove cues that rebuild the old state and add cues that make the
        target state easier to enter.
  \item \textbf{History shaping.} Target: emergence and path dependence.
        Evidence lane: medium + design inference.
        Use: repeat the same entry move in the same context so local history
        accumulates.
  \item \textbf{Minimum viable entry.} Target: barrier crossing under limited
        will.
        Evidence lane: medium + design inference.
        Use: define the smallest real action that still counts as entering the
        target regime.
  \item \textbf{External scaffolding.} Target: effective control input.
        Evidence lane: medium + design inference.
        Use: use schedules, partners, or deadlines when private control is too
        noisy to trigger reliable entry.
\end{itemize}
